\begin{beginnerstatement}[Kurz erklärt: Voraussetzungen und Installation]

\texttt{doctor} prüft lesend, ob die für das gewählte Profil wichtigen
Laufzeiten, Einstellungen, Pfade, Abhängigkeiten und Ports bereitstehen. Eine
Installation richtet benötigte Abhängigkeiten ein, also zusätzliche Programme
und Bibliotheken, auf denen das Projekt aufbaut. Ein \emph{Package Manager}
automatisiert das Herunterladen und Einrichten solcher Pakete. Node.js führt
das Frontend-Tooling aus; npm verwaltet dessen Pakete, und \texttt{npm ci}
installiert die im Lockfile festgelegten Stände. Python führt Backend und
Tooling aus; eine virtuelle Umgebung (\emph{venv}) trennt deren Pakete vom
übrigen System, während \texttt{uv} oder ersatzweise \texttt{pip} sie
installiert. Im v1.0.0-Kandidatenstand richtet \texttt{install} für jedes Profil
eine eigene \path{tools/.venv} mit Qualitätswerkzeugen wie Ruff ein; eine
zusätzliche \path{backend/.venv} enthält nur die Abhängigkeiten der Anwendung.
Git verwaltet Quellcodeänderungen und wird zum Klonen des
Repositories benötigt. Rust ist für die native Tauri-Schicht der
Desktopprofile erforderlich. Chromium ist der Browser für konfigurierte
Playwright-Tests und wird nur bei Bedarf eingerichtet; auch andere optionale
Werkzeuge werden profilabhängig behandelt.

\end{beginnerstatement}
